\documentclass[12pt,a4paper]{article}
\usepackage[T1]{fontenc}
\usepackage[width=2.5cm, height=2.5cm, left=2.5cm, right=2.5cm, top=2.5cm, bottom=2.5cm]{geometry}
\usepackage{graphicx}
\usepackage{mathtools}
\usepackage{amssymb}
\usepackage{amsthm}
\usepackage{thmtools}
\usepackage{tikz}
\usepackage{xcolor}
\usepackage{tcolorbox}
\usepackage{nameref}
\usepackage{hyperref}
\title{TP}
\author{KOIUTESSI Emmanuella }
\begin{document}
\vspace*{-3cm}
\noindent \includegraphics[width=0.25\textwidth]{unstim}
\hfill \includegraphics[width=0.2\textwidth]{ens}
\vspace{-2cm}
\begin{center}
	\includegraphics[width=0.3\textwidth]{benin}
\end{center}
\begin{center}
	RÉPUBLIQUE DU BENIN \\ MINISTÈRE DE L'ENSEIGNEMENT SUPÉRIEURE ET DE LA RECHERCHE SCIENTIFIQUE\\  UNIVERSITÉ NATIONALE DES SCIENCES TECHNOLOGIE INGÉNIERIE ET MATHÉMATIQUES \\ ÉCOLE NORMALE SUPÉRIEURE DE NATITINGOU 
\end{center}
\textcolor{blue}{\rule{18cm}{2pt}} 
\vspace*{2cm}
{\Large \textbf{I. MATHÉMATIQUES}} \\ 
Soit $f$ la fonction continue périodique. On considère ses coefficients : \\
$C_{n}(f)=\dfrac{1}{2 \pi} \displaystyle \int_{0}^{2 \pi}f(t)e^{-int}dt$
\begin{enumerate}
	\item Prouve que $\sum_{-\infty}^{+\infty}|C_{n}(f)|^2=\dfrac{1}{2 \pi} \displaystyle \int_{0}^{2 \pi}|f(t)|^2dt$
\end{enumerate}
{\Large {\textbf{II. PHYSIQUE (7pts)}}} \\
\textcolor{red}{\rule{17cm}{2.pt}} \\ \\

Démontrer l'invariant de la pseudo-norme du quadrivecteur impulsion : \\


$\left (\dfrac{E}{c} \right)^2- \left (\overrightarrow{p} \right)^2=m^2c^2$ \\

{\Large {\textbf{III. CHIMIE (6pts)}}} \\

\textcolor{green}{\rule{17cm}{2.pt}} \\ 
Calculer l'enthalpie libre de réaction 
 en fonction du quotient de réaction Q :\\
 
 $\Delta_{r}G$=$\Delta_{r}G°$+RT ln(Q)\\
 
 \colorbox{yellow}{\parbox{\linewidth}{En Afrique subsaharienne , la transformation digitale en santé se déploie dans un contexte singulier. Les systèmes de santé y sont caractériser par une insuffisance chronique de ressources financières , humaines et matérielles (Naabyonga-Orem \& Tachobya, 2016). Les infrastructures numériques , bien que croissantes avec l'essor de la téléphonie mobile et d' internet , restent inégales et parfois précaires . Les coupures d'électricité fréquentes , l'instabilité  des réseaux et la faible capacité technique du personnel    
 		constituent des contraintes majeures (Dovlo,2018). Cependant, ces fragilités  coexistes avec un réel potentiel }}

\end{document}
